\documentclass[landscape]{article}
\usepackage[pdftex]{graphicx}
\usepackage[utf8]{inputenc}
\usepackage{enumerate}
\usepackage{icomma}
\usepackage{amssymb}
\usepackage{tikz}
\usepackage{href-ul}
\hypersetup{
	colorlinks=true,
	linkcolor=blue,
	urlcolor=blue}
\usepackage{geometry}
\geometry{a4paper, top=15mm, left=15mm, right=15mm, bottom=15mm,
	headsep=10mm, footskip=12mm}

\begin{document}
	\begin{tikzpicture}(10,3)
		\draw[ultra thick](2,0) --(10,0) -- (10,1.3) --(2,1.3) -- (2,0);
		\draw[fill=black](2,0)-- (0,.6) -- (2,.6) -- (2,0);
		\node at (6,.6) {\large Lösungsblatt von \href{https://www.okuyakl.de}{www.okuyakl.de}};
	\end{tikzpicture}
	\vspace{0.5 cm}
	
{\bf Grenzwerte gebrochen rationaler Funktionen I -- Verhalten im Unendlichen}
\vspace{0.5 cm}

\renewcommand{\arraystretch}{3}
\begin{tabular}{|c|c|c|c|c|c|c|c|}
	\hline
	& Funktionsterm & Zählergrad & Nennergrad & Fall & $\lim\limits_{x\to\pm\infty}$ & Polynomdivision liefert & Asymptote / Näherungsfunktion \\
	\hline
a) & ${2x^2 -1 \over x+3}$ & Z=2 & N=1 & Z=N+1 & $\pm\infty$ & $f(x)=2x-6+{17\over x+3}$ & $y=2x-6$ \\
\hline
b) & ${x+3 \over x^2 -1}$ &Z=1 & N=2 & Z$<$N & 0 & -- & $y=0$\\
\hline
c) & ${x^3 +4x\over 2x^2-10x}$ & Z=3 & N=2 & Z=N+1 & $\pm\infty$ & $f(x)=0,5x + 2,5 + {29 \over 2x-10}$ & $y=0,5x + 2,5$\\
\hline
d) & ${x^4-16 \over x^2-2x}$& Z=4 & N=2 & Z$>$N+1 &$\infty$ & $f(x)=x^2 +2x + 4 + {8\over x}$ & $y=x^2 + 2x + 4 $ \\
\hline
e) & ${-0,1x^3 \over x^2+2x+1}$ & Z=3 &N=2 &Z=N+1 &$\mp \infty$ & $f(x)=-0,1x+0,2+{-0,3x-0,2 \over x^2+2x+1}$&  $y=-0,1x+0,2$\\
\hline
f) & ${x(x+5) \over 3x^2 }$ &Z=2 &N=2&Z=N & ${1 \over 3}$ & $f(x)={1\over 3}-{5\over 3x}$  &  $y={1 \over 3}$\\
\hline
g) & ${x^2-3x+2 \over x^3-2x^2 }$ &Z=2 &N=3 &Z$<$N & 0 & -- & $y=0$\\
\hline
h) & ${tx \over x+1 }$ &Z=1 &N=1 &Z=N & $t$  & $f(x)= t -{t \over x+1}$ &  $y=t$\\
\hline
i) &${3x-1 \over x-1}$  & Z=1 & N=1 & Z=N & 3 & $f(x)= 3 + {2\over x-1}$ &  $y=3$\\
\hline
j) & ${x^2-2 \over 2x-4}$ & Z=2 &N=1 & Z=N+1 & $\pm \infty$ & $f(x)= 0,5x + 1 + {1\over x-2}$ &  $y=0,5x + 1 $\\
\hline
\end{tabular}
\begin{center}
	\includegraphics[width=7 cm]{../../viecher/endcomic.pdf}
	
	Hier geht es zurück zum \href{https://www.okuyakl.de/math/m11ligrL070/aa070.pdf}{Aufgabenblatt}
\end{center}

\end{document}

